% =============================================================
% Prediction-horizon necessity study -- does a longer MPC lookahead
% improve closed-loop vessel-navigation tracking because it enables
% anticipatory control, or merely correlate with it? Standalone
% compilable excerpt, companion to methodology.tex and
% hardware_validation_section5.tex (\S\ref{sec:hw-validation}):
% those documents fix the Jacobian source as the independent
% variable; this one fixes the Jacobian source (the hybrid
% schedule) and varies ONLY the prediction horizon N, isolating
% horizon length as a causal variable in its own right. Written
% 2026-10-23 from a completed five-layer analysis (n=10/10 hardware
% campaign, a matched-configuration hard gate, a same-state
% command-divergence replay, and a correctly-recalibrated
% common-horizon nonlinear-oracle rollout).
%
% Layout note (2026-10-23, second revision): switched from
% \documentclass[journal]{IEEEtran} (two-column, ~245pt column
% width) to a single-column `article` layout with generous margins.
% Two successive attempts to fit this content into IEEEtran's
% narrow column (first via plain tabular/boxed, then via table*
% floats + stacked equations) both still overlapped -- a two-column
% journal layout is the wrong target for a dense, table- and
% equation-heavy methods appendix like this one; merge into the
% two-column paper body as prose only, once accepted, rather than
% re-fighting the column width here.
% =============================================================
\documentclass[11pt]{article}

\usepackage[margin=1in]{geometry}
\usepackage{amsmath}
\usepackage{amssymb}
\usepackage{amsfonts}
\usepackage{booktabs}
\usepackage{array}

\begin{document}

\section{Prediction-Horizon Necessity Study: Does Longer MPC Lookahead Improve Tracking Through Anticipation?}
\label{sec:horizon-necessity}

\subsection{Motivation and the Central Hypothesis}
\label{sec:ph-motivation}

The MPC controller of \S\ref{sec:control}\footnote{Cross-reference
into the companion methodology document; section label kept symbolic
here since the files are not yet merged.} re-solves a finite-horizon
quadratic program every tick, with prediction horizon $N$ a free
design parameter. The obvious experiment --- run the identical
controller at $N{=}15$ and $N{=}1$ and compare tracking error ---
answers only a weak, correlational question. The claim this section
is designed to support or refute is the stronger, mechanistic one:
\begin{equation}
\boxed{
N{=}15 \;\rightarrow\; \text{better anticipatory action}
\;\rightarrow\; \text{smaller subsequent tracking error}
}
\label{eq:causal-chain}
\end{equation}
not merely
\begin{equation}
\boxed{\mathrm{RMSE}_{15} < \mathrm{RMSE}_{1}.}
\label{eq:weak-claim}
\end{equation}
Equation~\eqref{eq:weak-claim} is necessary but not sufficient
evidence for \eqref{eq:causal-chain}: a horizon-dependent timing
confound, a trivial constant-offset effect, or an artifact of which
states each condition happens to visit could all produce
\eqref{eq:weak-claim} without any genuine anticipatory mechanism. The
analysis below is organized into five layers, each designed to rule
out one specific alternative explanation before the next layer's
result is trusted.

\subsection{Experimental Design: Five Layers}
\label{sec:ph-design}

All conditions share the dipole-aligned floor240 vessel plan, the
contact-model hybrid Jacobian schedule (contact-model joint columns,
no-contact-model insertion column), and every controller weight,
constraint, and rate limit of \S\ref{sec:control}; only the
prediction horizon $N\in\{1,15\}$ differs. Ten clean hardware
repetitions per condition were collected (two further conditions on
the no-dipole-alignment variant of the same plan replicate the same
pattern and are reported in \S\ref{sec:ph-nodipole} for completeness).

\begin{enumerate}
\item \textbf{Layer 1 (matched-comparison hard gate).} Every logged
  controller-configuration field is diffed between the two
  conditions; real-time control-loop cadence and QP solver health are
  compared explicitly, since a horizon-dependent solve-time increase
  could silently change the realized sample period.
\item \textbf{Layer 2 (run-level physical tracking).} Per-run RMSE, median,
  $p_{90}$, $p_{95}$, $\max$ of the tracking error $e(s)$, parametrized
  by physical path progress $s$ (not tick index), plus lag and
  cross-track decomposition. Run-level (never tick-level) bootstrap
  confidence interval and an exact permutation test on
  $\Delta_{\mathrm{RMSE}} = \mathrm{median}_r\,\mathrm{RMSE}_r^{(15)} -
  \mathrm{median}_r\,\mathrm{RMSE}_r^{(1)}$.
\item \textbf{Layer 3 (region-C signature, inconclusive by design).} Tests
  whether command divergence precedes error divergence specifically
  around the established contact-sensitive region C
  ($s\in[55,65]\,\mathrm{mm}$). Found uninformative here --- see
  \S\ref{sec:ph-layer3}.
\item \textbf{Layer 4 (same-state command replay).} At matched anchor
  states $x_k$ drawn from the real hardware logs, two controller
  objects --- identical in every respect except $N$ --- are each
  solved once, holding $x_k$, the delay-filter residual, the previous
  command, the Jacobian schedule, and all weights fixed. Only the
  first command $u_0^{(N)}(x_k)$ is compared.
\item \textbf{Layer 5 (consequence of the divergence).} Two independent
  tests link the Layer-4 command divergence to subsequent physical
  outcome: a delay-corrected correlation against the real hardware's
  own future error (\S\ref{sec:ph-5a}), and a common-horizon nonlinear
  oracle rollout (\S\ref{sec:ph-5b}) that evaluates both candidate
  first commands under an identical future policy and an identical,
  correctly-calibrated contact-model ground truth.
\end{enumerate}

\subsection{Layer 1: The Comparison Is Matched, With One Honest Caveat}
\label{sec:ph-layer1}

Every logged field in \texttt{summary.json}/\texttt{controller\_metadata.json}
is identical between the $N{=}15$ and $N{=}1$ conditions except
\texttt{horizon} itself (and the derived \texttt{run\_name}). No
Jacobian-source, weight, delay, or constraint difference survives this
check.

A genuine timing confound does exist and is reported rather than
hidden (Table~\ref{tab:timing}). $N{=}15$ runs its control loop
$18.7\%$ slower in real time than $N{=}1$. Critically, this bias runs
\emph{against} the hypothesis under test: a feedback loop that samples
state less often should track worse, not better. If $N{=}15$
nonetheless wins decisively (\S\ref{sec:ph-layer2}) despite this
handicap, the result is \emph{more}, not less, credible.
Commanded-vs-achieved velocity realization is statistically identical
between conditions (servo RMS error $0.00569$ vs.\
$0.00575\,\mathrm{rad}$), ruling out an actuation-fidelity explanation
for the slower loop.

\begin{table}[h]
\centering
\caption{Layer 1: realized control-loop timing, $N{=}15$ vs.\ $N{=}1$ (nominal tick period $0.100\,\mathrm{s}$ in both).}
\label{tab:timing}
\begin{tabular}{@{}lcc@{}}
\toprule
 & $N{=}15$ & $N{=}1$ \\
\midrule
Mean solver time & $\sim$20--23\,ms & 4.46\,ms \\
QP hits its own time limit & 150/5513 (2.7\%) & 0/6367 (0\%) \\
Realized tick period & 0.1253\,s & 0.1056\,s \\
\bottomrule
\end{tabular}
\end{table}

\subsection{Layer 2: Complete, Non-Overlapping RMSE Separation}
\label{sec:ph-layer2}

\begin{table}[h]
\centering
\caption{Layer 2: run-level physical tracking, $n{=}10$ repetitions per condition.}
\label{tab:rmse}
\begin{tabular}{@{}lcc@{}}
\toprule
 & $N{=}15$ & $N{=}1$ \\
\midrule
Median RMSE & $0.708\,\mathrm{mm}$ & $2.232\,\mathrm{mm}$ \\
Range & $[0.626,0.786]\,\mathrm{mm}$ & $[2.035,2.483]\,\mathrm{mm}$ \\
Mean lag RMS & $0.535\,\mathrm{mm}$ & $1.122\,\mathrm{mm}$ \\
Mean cross-track RMS & $0.451\,\mathrm{mm}$ & $1.879\,\mathrm{mm}$ \\
\bottomrule
\end{tabular}
\end{table}

The two ten-run samples do not overlap at all: the worst $N{=}15$ run
($0.786\,\mathrm{mm}$) beats the best $N{=}1$ run
($2.035\,\mathrm{mm}$). The run-level bootstrap 95\% confidence
interval on $\Delta_{\mathrm{RMSE}}$ is $[-1.653,-1.417]\,\mathrm{mm}$,
entirely negative. The exact permutation test over all
$\binom{20}{10}=184{,}756$ possible run-label assignments gives
$p=0.00076$ (only 140 assignments produce a gap this large or
larger). Cross-track error degrades roughly $4\times$ going from
$N{=}15$ to $N{=}1$, versus only $2\times$ for lag --- $N{=}1$'s
deficiency is disproportionately lateral, the signature expected if
it cannot anticipate upcoming steering demand.

Equation~\eqref{eq:weak-claim} is therefore established beyond
reasonable doubt. The remaining layers test whether
\eqref{eq:causal-chain} is too.

\subsection{Layer 3: An Honest Null, Not a Negative}
\label{sec:ph-layer3}

A ``does command divergence precede error divergence'' test, binned
by path progress $s$, found both $e(s)$ and the command signals
$\|u(s)\|,\|\Delta u(s)\|$ already significantly separated between
conditions at the very first bin ($s=1.5\,\mathrm{mm}$) --- a direct
consequence of Layer 2's complete separation having no overlap
anywhere along the path. The test is uninformative by construction
once the two conditions never agree at any point, not evidence
against anticipation. This motivates Layers 4--5, which compare
counterfactual commands \emph{at the same state} rather than two
already-diverged trajectories.

\subsection{Layer 4: Horizon Length Materially Changes the Optimized Action}
\label{sec:ph-layer4}

At $n=3955$ same-state anchors (full 10+10 dataset), holding
$x_k,u_{k-1}$, the delay-filter residual, the Jacobian schedule, and
all weights fixed and varying only $N$, the mean angular divergence
between the two candidate first commands,
$\theta_u = \angle(u_0^{(15)},u_0^{(1)})$, is shown in
Table~\ref{tab:theta}.

\begin{table}[h]
\centering
\caption{Layer 4: same-state command divergence $\theta_u$ by path region (pre-region $s{\in}[40,55]\,\mathrm{mm}$, region C $s{\in}[55,65]\,\mathrm{mm}$).}
\label{tab:theta}
\begin{tabular}{@{}lccc@{}}
\toprule
 & All anchors & Pre-region & Region C \\
\midrule
Mean $\theta_u$ & $35.53^{\circ}$ & $34.20^{\circ}$ & $41.06^{\circ}$ \\
\bottomrule
\end{tabular}
\end{table}

Commands are never close to identical ($\theta_u$ never approaches
$0^{\circ}$ anywhere in the dataset), ruling out the null
``$u_0^{(15)}\approx u_0^{(1)}$ everywhere, so horizon length is not
explaining the hardware difference.'' The $+20\%$ step-up entering
region C is consistent with --- though on its own insufficient to
prove --- the divergence concentrating where lookahead should matter
most.

\subsection{Layer 5a: A Confounded Correlation, Correctly Discarded}
\label{sec:ph-5a}

Defining a delay-corrected future-error statistic
$E_{\mathrm{future}}(k) = \tfrac{1}{H_f}\sum_{h=d}^{d+H_f-1}\|e_{k+h}\|$
($d{=}2$ delay samples, $H_f{=}10$) and correlating it against the
Layer-4 divergence magnitude within each condition's own real
trajectory gives
\begin{equation}
\mathrm{Spearman}(\theta_u, E_{\mathrm{future}}) = +0.287\ (N{=}15),
\quad +0.700\ (N{=}1),
\end{equation}
both strongly positive and highly significant ($p\ll 0.001$) --- the
\emph{wrong} sign for a naive ``bigger departure causes better
outcome'' reading. This is explained by a common-cause confound,
not a refutation: harder states independently produce both larger
same-state disagreement and worse subsequent error in both
conditions, without one causing the other. This result is correctly
treated as non-identifying, not as evidence against $N{=}15$, and
motivates the counterfactual design of Layer 5b.

\subsection{Layer 5b: Common-Horizon Oracle --- the Decisive Test}
\label{sec:ph-5b}

At each anchor, both candidate first commands $u_0^{(15)}(x_k)$ and
$u_0^{(1)}(x_k)$ are rolled forward under an \emph{identical} policy
(the candidate for one step, then the plan's own reference velocity
for the remaining $H_c{-}1{=}7$ steps) and scored against an
\emph{identical}, correctly-calibrated nonlinear contact-model oracle
evaluated against the plan's true target tip position (not a
joint-state proxy --- see \S\ref{sec:ph-oracle-bug} for why that
distinction mattered). This isolates the horizon's causal contribution
from state-difficulty confounds by construction: both candidates face
the same future, the same ground truth, and only their first action
differs.

\begin{table}[h]
\centering
\caption{Layer 5b: common-horizon oracle outcome. $\Delta_{\mathrm{rms}} = \mathrm{rms}(u_0^{(1)}) - \mathrm{rms}(u_0^{(15)})$; $\Delta_{\mathrm{rms}}{>}0$ means the $N{=}15$ command has lower predicted cost.}
\label{tab:oracle}
\begin{tabular}{@{}lccc@{}}
\toprule
 & $n$ & $P(\Delta_{\mathrm{rms}}{>}0)$ & Median $\Delta_{\mathrm{rms}}$ \\
\midrule
Overall & 334 & $0.617$ & $+0.0091\,\mathrm{mm}$ \\
Pre-region & 15 & $0.600$ & $+0.0031\,\mathrm{mm}$ \\
Region C & 264 & $\mathbf{0.636}$ & $+0.0116\,\mathrm{mm}$ \\
Elsewhere & 55 & $\mathbf{0.527}$ ($\approx$chance) & $+0.0018\,\mathrm{mm}$ \\
\bottomrule
\end{tabular}
\end{table}

The single most mechanistically informative number in
Table~\ref{tab:oracle} is the region contrast: the advantage is
\emph{real and elevated specifically inside region C} (64\%) and
\emph{collapses to chance everywhere else} (53\%). This is the
anticipation signature \eqref{eq:causal-chain} predicts --- not a
uniform background advantage that happens to coincide with better
overall tracking, but one that concentrates precisely where the
path's geometric demand is highest and lookahead has the most to
offer.

\subsubsection{A Methodological Pitfall That Nearly Produced a False Null}
\label{sec:ph-oracle-bug}

An initial version of this test scored both candidates against
$\mathrm{oracle}(\mathrm{reference.state}[\mathrm{idx}])$ --- the
offline plan's own joint-space solution at the matching index --- as
ground truth. This produced absurd magnitudes (oracle RMS of
$9$--$19\,\mathrm{mm}$, later traced to as much as $50\,\mathrm{mm}$),
physically impossible given the vessel lumen's own
$5.6$--$10.2\,\mathrm{mm}$ radius, and a result trending \emph{toward}
chance as more data arrived ($0.565\to0.548$) rather than away from
it. Two independent, compounding errors were found and fixed:
\begin{enumerate}
\item The oracle adapter was built without an explicit
  \texttt{initial\_poses} argument, silently falling back to a
  generic default beam-base pivot ($z{=}{-}0.0166\,\mathrm{m}$) instead
  of this plan's real, recalibrated pivot
  ($z{=}{-}0.0396\,\mathrm{m}$) --- the identical $23\,\mathrm{mm}$
  bug class this project had already fixed once before, reintroduced
  by omission.
\item Even with the pivot corrected, scoring against
  $\mathrm{oracle}(\mathrm{reference.state}[\mathrm{idx}])$ remained
  invalid: this controller runs with a reference-input nulling option
  in its cost (\texttt{diagnostic\_zero\_input\_reference\_in\_R}) that
  deliberately lets the arm's kinematically redundant joint drift away
  from the plan's own joint-space solution while preserving tip
  tracking. The true nonlinear model is \emph{not} fully agnostic to
  that redundant motion the way the live controller's own simplified
  Jacobian is, so the reference joint state silently drifts out of
  registration with physical reality as the run progresses (an
  oracle-vs-oracle discrepancy growing from $0.07\,\mathrm{mm}$ at
  tick 0 to $4.4\,\mathrm{mm}$ by tick 365, tracking the growing
  joint-space divergence almost exactly) --- while real,
  vision-measured tracking error stayed under $1\,\mathrm{mm}$
  throughout. The fix was to score against the plan's actual target
  tip position (arc-length indexed, invariant to joint
  configuration), eliminating both the frame-calibration bug and the
  redundant-DOF confound simultaneously, and as a side effect halving
  the oracle evaluations required.
\end{enumerate}
After both fixes, oracle-predicted tip position matched real
vision-measured tip position to within $2$--$6\,\mathrm{mm}$ (a
physically sensible residual, itself growing mildly with the same
redundant-DOF effect), and the result above emerged. This is recorded
here because it is exactly the kind of silent, large, and
state-dependent methodological error that a weaker check (e.g.\
trusting internal self-consistency between two oracle evaluations,
rather than validating against an independent ground truth) would not
have caught.

\subsection{Synthesis: What Causes the Improvement}
\label{sec:ph-synthesis}

The five-layer chain supports a specific, qualified version of
\eqref{eq:causal-chain}:
\begin{enumerate}
\item The configuration is genuinely matched (\S\ref{sec:ph-layer1}), and
  the one real confound (slower realized loop rate under $N{=}15$)
  biases against, not for, the observed result.
\item The hardware effect is large, decisive, and non-overlapping
  across ten independent repetitions per condition
  (\S\ref{sec:ph-layer2}).
\item Horizon length is not a nominal parameter: it changes the
  optimized command by a large margin ($\sim\!35^{\circ}$) at
  identical states (\S\ref{sec:ph-layer4}).
\item That command divergence corresponds to a genuinely lower
  predicted future cost under a controlled, common-policy,
  correctly-calibrated nonlinear rollout (\S\ref{sec:ph-5b}), and the
  size of that advantage is itself concentrated in the one region
  where anticipation should matter most (64\% in region C vs.\ 53\%,
  indistinguishable from chance, elsewhere).
\end{enumerate}
The mechanism is that a 15-step ($1.5\,\mathrm{s}$) horizon lets the
optimizer account for the upcoming high-curvature, contact-sensitive
segment of the path before the tracking error from it has appeared,
while a 1-step horizon is purely reactive and can only respond once
error is already present; this is directly visible as cross-track
error (the lateral, anticipation-sensitive component) degrading
roughly twice as much as lag error (the along-path, purely reactive
component) when the horizon is shortened (\S\ref{sec:ph-layer2}).

The appropriate strength of claim is, however, a qualified one. The
Layer 5b effect size (62\% overall, 64\% in region C) is real,
region-localized, and statistically distinguishable from the
chance-level 53\% observed outside that region, but it is moderate,
not overwhelming. The $3.2\times$ hardware RMSE gap is far larger than
a single-tick 62\%-vs-chance advantage can plausibly explain in one
step; the most defensible reading is that anticipation provides a
real, demonstrable, region-localized edge at the level of individual
decisions, and that this edge compounds across hundreds of ticks into
the much larger aggregate advantage observed on hardware ---
distinguishing it from a competing (and not fully excluded) hypothesis
that most of the aggregate gap instead reflects the horizon-$1$
controller's reactive-only commands repeatedly overshooting and
correcting along the way, a dynamic whose cumulative cost could exceed
what any single-tick oracle comparison captures.

\subsection{No-Dipole-Alignment Replication}
\label{sec:ph-nodipole}

The identical $N\in\{1,15\}$ comparison, repeated on the no-dipole-
alignment variant of the same plan (hybrid schedule, otherwise
identical configuration), reproduces the Layer 2 pattern closely:
median RMSE $0.66$--$0.68\,\mathrm{mm}$ at $N{=}15$ versus
$1.75$--$1.88\,\mathrm{mm}$ at $N{=}1$ (complete separation, $n{=}10$
each), confirming the horizon effect is not an artifact specific to
the dipole-alignment cost term.

\end{document}
